%SOP Template 
% Version 01 Added revision date
% Version 03 Added TOC and acknowledgements
%           New SOP3_alpha.cls


\documentclass[12pt]{SOP5_beta}\usepackage[]{graphicx}\usepackage[]{xcolor}
% maxwidth is the original width if it is less than linewidth
% otherwise use linewidth (to make sure the graphics do not exceed the margin)
\makeatletter
\def\maxwidth{ %
  \ifdim\Gin@nat@width>\linewidth
    \linewidth
  \else
    \Gin@nat@width
  \fi
}
\makeatother

\definecolor{fgcolor}{rgb}{0.345, 0.345, 0.345}
\newcommand{\hlnum}[1]{\textcolor[rgb]{0.686,0.059,0.569}{#1}}%
\newcommand{\hlsng}[1]{\textcolor[rgb]{0.192,0.494,0.8}{#1}}%
\newcommand{\hlcom}[1]{\textcolor[rgb]{0.678,0.584,0.686}{\textit{#1}}}%
\newcommand{\hlopt}[1]{\textcolor[rgb]{0,0,0}{#1}}%
\newcommand{\hldef}[1]{\textcolor[rgb]{0.345,0.345,0.345}{#1}}%
\newcommand{\hlkwa}[1]{\textcolor[rgb]{0.161,0.373,0.58}{\textbf{#1}}}%
\newcommand{\hlkwb}[1]{\textcolor[rgb]{0.69,0.353,0.396}{#1}}%
\newcommand{\hlkwc}[1]{\textcolor[rgb]{0.333,0.667,0.333}{#1}}%
\newcommand{\hlkwd}[1]{\textcolor[rgb]{0.737,0.353,0.396}{\textbf{#1}}}%
\let\hlipl\hlkwb

\usepackage{framed}
\makeatletter
\newenvironment{kframe}{%
 \def\at@end@of@kframe{}%
 \ifinner\ifhmode%
  \def\at@end@of@kframe{\end{minipage}}%
  \begin{minipage}{\columnwidth}%
 \fi\fi%
 \def\FrameCommand##1{\hskip\@totalleftmargin \hskip-\fboxsep
 \colorbox{shadecolor}{##1}\hskip-\fboxsep
     % There is no \\@totalrightmargin, so:
     \hskip-\linewidth \hskip-\@totalleftmargin \hskip\columnwidth}%
 \MakeFramed {\advance\hsize-\width
   \@totalleftmargin\z@ \linewidth\hsize
   \@setminipage}}%
 {\par\unskip\endMakeFramed%
 \at@end@of@kframe}
\makeatother

\definecolor{shadecolor}{rgb}{.97, .97, .97}
\definecolor{messagecolor}{rgb}{0, 0, 0}
\definecolor{warningcolor}{rgb}{1, 0, 1}
\definecolor{errorcolor}{rgb}{1, 0, 0}
\newenvironment{knitrout}{}{} % an empty environment to be redefined in TeX

\usepackage{alltt}
\usepackage[english]{babel}
% \usepackage{blindtext}
% \usepackage{lipsum}

%\documentclass{article}

%\documentclass[12pt]{~/github/SOPs/SOP_Template/SOP}

\title{DNA Isolation from Soil -- DNeasy PowerSoil ProKit}
\date{7/02/2025}
\author{Ghena Kubba}
\approved{TBD}
\ReviseDate{\today}
\SOPno{39 v.01}
\IfFileExists{upquote.sty}{\usepackage{upquote}}{}
\begin{document}
\maketitle

\section{DNA Extraction }

\NP This procedure is to extract DNA from soils using a QIAGEN DNeasy Powersoil Pro Kit. 
%need trademark here

\NP The applications of this SOP include ribsomal DNA, community profiling and PCR/qPCR for specific DNA sequences from soil microbes. 

\section{Summary of Method}

Soil samples are homogenized by mechanical treatment or collected in a manner that does not require additional treatment (i.e. samples suspended in water/solvent). 

The PowerSoil\textsuperscript{\textregistered} DNA Isolation Kit is effective at removing PCR inhibitors from even the most difficult soil types. This kit features a novel bead tube and optimized chemistry for a more efficient lysis of soil bacteria and fungus. The kit is shown to be the most effective in isolating high-quality DNA

\tableofcontents

\newpage

\section{Acknowledgements}

\NP This SOP template is based off.........
was originally written by Aparna... to extract DNA from surface waters. Initial test of the methods were done by Alejandro Guerrero in the summer of 2016. 

\section{Definitions}

\NP
\textbf{Supernatant:} The clear liquid that lies above a solid residue after centrifugation. In this protocol, the supernatant contains the DNA and other soluble components separated from cell debris and particulate matter
\NP
\textbf{Pellet:} The dense solid material that collects at the bottom of a tube after centrifugation. In this protocol, the pellet typically consists of soil particles, cell debris, and other insoluble contaminants that are separated from the DNA-containing solution.
\NP
\textbf{Lysis:} The process of breaking open cells to release their internal contents, including DNA. In this protocol, lysis is achieved through mechanical disruption (bead beating) and chemical reagents that break down cellular structures.
\NP
\textbf{Elution:} The process of releasing purified DNA from the silica membrane of the spin column into a clean microcentrifuge tube. After DNA binds to the membrane during earlier steps, a low-salt buffer (usually Solution C6) is added to the column and centrifuged. This buffer disrupts the interaction between the DNA and the membrane, allowing the DNA to be collected in the final eluate. The resulting solution contains the purified DNA, ready for downstream applications such as qPCR, ddPCR, or sequencing.
\NP
\textbf{Binding Buffer:} A solution that creates optimal conditions for DNA to bind to the silica membrane in the spin column. This is typically done by adjusting salt concentration and pH.
\NP
\textbf{Wash Buffer:} A solution used to remove residual contaminants such as proteins, salts, and inhibitors while the DNA remains bound to the membrane. Multiple washes ensure high DNA purity.
\NP
\textbf{Elution Buffer:}A low-salt buffer provided in the kit that is used to release the bound DNA from the silica membrane into a final elution tube. It ensures the DNA is in a suitable solution for downstream applications like PCR or sequencing.
\NP
\textbf{Centrifugation:} The process of spinning samples at high speed to separate components based on density. In this protocol, centrifugation is used to isolate supernatant from pellet and to move liquids through the spin column.

\section{Biases and Interferences}
\NP The extraction process is supposed to isolate DNA from other cellular or other organic materials. However, certain compounds commonly present in soil—such as humic acids, polyphenols, and residual proteins—can co-purify with DNA and interfere with downstream analyses.

\NP When meausuring the DNA yield, these compounds may interfere with the measurements and thus, suggest you have more DNA that was extracted in reality. 

\NP Mesuring the absorbance ratios of several wavelength gives a reasonble estimate of DNA purity. See the QA/QC section to measure these contaminants.

\begin{description}
  \item[A260/A280 Ratio] Nucleic acids, DNA and RNA, absorb at 260nm. For a pure sample, a well defined peak (no shoulders or wiggles) at 260nm is expected.Several factors, however, can influence the accuracy of the 260/280 and 260/230 ratios. Readings from very dilute samples will have very little difference between the absorbance at 260 and 280nm leading to inaccurate ratios.  The type(s) of protein present will also have an effect.  Absorbance in the UV range by proteins is primarily the result of aromatic ring structures. Phenol and other contaminants can also absorb at 280 nm and can affect the ratio calculation. Phenol absorbs with a peak at 270nm. \textbf{Nucleic acid preparations uncontaminated by phenol should have an 260/280 ratio of around 1.8. } The pH of the solution can also affect the 260/280 ratio, with acidic solutions having a lower ratio of up to 0.2-0.3 and alkaline solutions having an increased ratio by a similar amount.
  \item[A260/A230 Ratio] Pure RNA has an A260/A280 ratio of 2.0, therefore if a DNA sample has an 260/280 ratio of greater than 1.8 this could suggest RNA contamination. The 260/280 ratio is a secondary measure of nucleic acid purity. This ratio for pure samples are often higher than the respective 260/280 ratio values. Strong absorbance around 230nm can indicate that organic compounds or chaotropic salts are present in the purified DNA.  A ratio of 260nm to 230nm can help evaluate the level of salt carryover in the purified DNA. The lower the ratio, the greater the amount of salt present. As a guideline, \textbf{the 260/230 ratio should be greater than 1.5, ideally close to 1.8.} Urea, EDTA, carbohydrates and phenolate ions all have absorbance near 230nm. A reading at 320nm will indicate if there is turbidity in the solution, another indication of possible contamination.  

\end{description}


\section{Health and Safety}

\NP The risks involved with using this kit are mainly related to the chemicals that are in use. As listed below, safety precautions should be taken at all times, but especially when handling hazardous reagants. While following safety protocol, it is advised that the materials and equipment are kept organized to avoid contamination and thus yield good results. 

\NP Avoid all skin contact with reagents in this kit. 

\NP In case of contact, wash thoroughly with water. Do not
ingest. See Safety Data Sheets for emergency procedures in case
of accidental ingestion or contact. 

\NP All SDS information is available upon request (800-362-7737) or on our web site at www.qiagen.com/safety. Reagents labeled flammable should be kept away from open flames and sparks.

\NP \textbf{WARNING:} Solution C5 and EA are flammable 
\NP \textbf{WARNING:} DO NOT add bleach or acidic solutions directly to the sample preparation waste. Solution CD1 and CD3 contain chaotropic salts, which can form highly reactive compounds when combined with bleach. If liquid containing these buffers is spilt, clean with a suitable labratory detergent and water. 
 
\subsection {Safety and Personnnel Protective Equipment}

\NP In addition, to appropriate laboratory attire, always wear gloves, laboratory coat, and protective googles when following this SOP.

\section{Personnel \& Training Responsibilities}

\NP Training is required before this SOP can be used. 

\NP Researchers using this SOP should be trained for the following SOPs:

\begin{itemize*}
  \item SOP01 Laboratory Safety
  \item SOP03 Handling of Hazardous Materials
  \item SOP09 Using Balances, Pippettes, and Glassware
  \item SOP10 Quality Control and Quality Assurance
  \item SOP14 Microcentrifuge
  \item SOP16 Using Laboratory Refrigerators and Freezers
\end{itemize*}

\section{Required Materials}

\subsection*{Equipment}

\NP Microcentrifuge with rotor capable of reaching 16,000 g

\NP Pipettor 50 $\mu$L - 750 $\mu$L

\NP Vortex-Genie 2

\NP Vortex Adapter for 24 (1.5-2 ml) tubes
%need trademark here

\NP (Optional) Bead-Beater: If a vortex is not available, or the sample requires stronger homogenization, the kit contains bead tubes suitable for high-powered bead beating and can be used in junction with the PowerLyzer® 24 Homogenizer (110/220V) (cat. no. 13155) or the TissueLyser II (cat. no. 85300) using a 2 ml Tube Holder Set (cat no. 11993).

\NP Sharpies

\NP Laboratory Notebook

\subsection*{Reagents}

\NP Reagents are provided by DNeasy PowerSoil Pro Kit 

\subsection*{Consumables}

\NP RNAase free gloves (cat \# ???)

\NP Disposable pippet tips (for 100 preps) 

\NP Vortex adapter for bead tubes

\begin{itemize*}
  \item 100 --- 2 $\mu$L tips
  \item 50 --- 200 $\mu$L tips
  \item 700 --- 1000 $\mu$L tips
  \item 50 --- 5000 $\mu$L tips

\end{itemize*}

\NP Weighing Paper

\NP Scupula, spatula, tweezers

\subsection*{Kit Contents}
\NP PowerBead Pro Tubes (50)
\NP MB Spin Columns (50)
\NP Solution CD1 (40 ml)
\NP Solution CD2 (15 ml)
\NP Solution CD3 (35 ml)
\NP Solution EA (36 ml)
\NP Solution C5 (30 ml)
\NP Solution C6 (9 ml)
\NP Microcentrifuge Tubes - 2 ml (100)
\NP Elution Tubes - 1.5 ml (50)
\NP Collection Tubes - 2 ml (100)
\NP Quick-Start Protocol (1)


\section{Estimated Time (batch of 12)}

\NP This procedure requires approximately 5-6 hours depending on the number of samples,and the time it takes to homogenize starting material and complete other preliminary steps.

\NP Batches of 12 are easiest to follow for new researchers and takes about 5 hours, while batches of 24 are more efficient, but can take more than 8 hours.

\section{Sample Collection, Preservation, and Storage}

\NP \textbf{Solution CD2 should be stored at 2-8\degree C upon arrival.} All other components and reagents of the DNeasy Powersoil Pro Kit can be stored at room temperature (15-25\degree C) until the expiration date printed on the label.

\NP Before the extraction process, soils should be either freshly collected or frozen... MORE here...

\NP After extraction, the DNA shall stored in the  –30 to –15\degree C or –90 to –65\degree C.

\section{Procedure}

\subsection*{Preliminary Steps} 
 
\NP Put soils in order to facilate an efficient weighing process. 

\NP Label PowerBead Pro Tubes (provided) in the same order as soils.

\NP Spend a few minutes homogenizing each soil before sampling.

\NP To the PowerBead Tubes, add 250 mg (0.25 grams) of soil sample.

\noindent What's happening: After your sample has been loaded into the PowerBead Tube, the next step is a homogenization and lysis procedure. The PowerBead
Tube contains a buffer that will (a) help disperse the soil particles, (b) begin to dissolve humic acids and (c) protect nucleic acids from degradation.

\NP Gently vortex to mix.

\noindent What's happening: Gentle vortexing mixes the components in the PowerBead Tube and begins to disperse the sample in the PowerBead Solution.

\NP Check Solution CD3. If Solution CD3 is precipitated, heat solution to 60 \degree C until the precipitate has dissolved before use.

\noindent What's happening: Solution CD3 contains SDS and other disruption agents required for complete cell lysis. In addition to aiding in cell lysis, SDS is an anionic detergent that breaks down fatty acids and lipids associated with the cell membrane of several organisms. If it gets cold, it will form a white precipitate in the bottle. Heating to 60 \degree C will dissolve the SDS and will not harm the SDS or the other disruption agents. Solution CD3 can be used while it is still warm.

\NP Spin the PowerBead Pro Tube briefly to ensure beads have settled at the bottom. In addition to the 250 mg of soil, also add 800 $\mu$L of Solution CD1. Vortex briefly to mix

\noindent What's Happening: After the sample has been loaded into the PowerBead Pro Tube, the next step is a homogenization and lysis procedure. The PowerBead Pro Tube contains a buffer that will (a) help disperse the soil particles, (b) begin to dissolve humic acids, and (c) protect nucleic acids from degradation. Gentle vortexing mixes the components in the PowerBead Pro Tube and begins to disperse the sample in the buffer.

\subsection*{Homogenization Process} 
\NP Homogenize samples thoroughly using one of the following methods: 
\begin{itemize*}
  \item Secure the PowerBead Pro Tube horizontally on a Vortex Adapter for 1.5–2 ml tubes (cat. no. 13000-V1-24). Vortex at maximum speed for 10 min. \textbf{Note:} If using Vortex Adapter for more than 12 preps simultaneously, increase the vortexing time by 5–10 min \textbf{Note:} Using the Vortex Adapter will maximize homogenization, which can lead to higher DNA yields. Avoid using tape, which can become loose and result in reduced homogenization efficiency, inconsistent results, and reduced yields.
  \item Use a PowerLyzer 24 Homogenizer. PowerBead Pro Tubes must be properly balanced in the tube holder of the PowerLyzer 24 Homogenizer. We recommend homogenizing the soil at 2000 rpm for 30 s, pausing for 30 s, then homogenizing again at 2000 rpm for 30 s. \textbf{Note:} Homogenizing samples at higher speeds (up to 4000 rpm) may increase yields but result in more fragmented DNA.
\end{itemize*}

\noindent The vortexing step is critical for complete homogenization and cell lysis. Cells are lysed by a combination of chemical agents from steps 1-4 and mechanical shaking introduced at this step. By randomly shaking the beads in the presence of disruption agents, collision of the beads with microbial cells will cause the cells to break open.

\subsection*{Inhibitor and Contaminant Removal}
\NP Centrifuge the PowerBead Pro Tube at 15,000 x g for 1 min.
\NP Transfer 500-600 $\mu$L of the supernatant to a clean 2 ml Microcentrifuge Tube (provided). Avoid disturbing the pellet. There may also still be soil particles.
\NP Add 200 $\mu$L of Solution CD2 and vortex for 5 s.
\NP Centrifuge at 15,000 x g for 1 minute
\NP Avoiding the pellet, transfer up to 700 $\mu$L of supernatant to a clean 2 ml tube. 
\NP Add 600 $\mu$L of Solution CD3 and vortex for 5 seconds
\noindent What's happening: CD2 and CD3 contain Inhibitor Removal Technology (IRT) reagents that precipitate proteins, humic acids, and other inhibitors that interfere with PCR and other downstream applications.
\subsection*{DNA Binding and Washing}
\NP Immediately after vortexing with CD3, load 650 $\mu$L of the lysate onto an MB Spin Column in a 2 mL Collection Tube.
\NP Centrifuge at 15,000 × g for 1 minute. Discard the flow-through.
Repeat until all of the lysate has been processed and passed through the MB Spin Column.
\NP Carefully place the MB Spin Column into a clean 2 ml Collection Tube (provided). Avoid splashing any flow-through onto the MB Spin Column.
\NP Add 500 $\mu$L of Solution EA to the MB Spin Column. Centrifuge at 15,000 x g for 1 min.
\NP Discard the flow-through and place the MB Spin Column back into the same 2 ml Collection Tube.
\NP Add 500 $\mu$L of Solution C5 to the MB Spin Column. Centrifuge at 15,000 x g for 1 min.
\NP Discard the flow-through and place the MB Spin Column into a new 2 ml Collection Tube (provided).
\NP Centrifuge at up to 16,000 x g for 2 min. Carefully place the MB Spin Column into a new 1.5 ml Elution Tube (provided). \textbf{Note:} This spin removes residual Solution C5. It is critical to remove all traces of
Solution C5 because the ethanol in it can interfere with downstream DNA applications, such as PCR, restriction digests, and gel electrophoresis.
\noindent What's Happening: Ethanol washes remove residual salts, proteins, and inhibitors while allowing DNA to stay bound. The final spin ensures no ethanol is left, which can inhibit downstream reactions.

\subsection*{Elution}
\NP Add 50-100 of $\mu$L of Solution C6 to the center of the while filter membrane. \textbf{Note:} Placing Solution C6 in the center of the small white membrane will make sure the entire membrane is wet. This will result in a more efficient and complete release of the DNA from the MB Spin Column filter membrane. As Solution C6 passes through the silica
membrane, DNA that was bound in the presence of high salt is selectively released by Solution C6 (10 mM Tris), which lacks salt.
\NP Centrifuge at 15,000 x g for 1 min. Discard the MB Spin Column. The DNA is now ready for downstream applications.
\NP The DNA should be stored frozen  (–30 to –15 \degree C or –90 to –65 \degree C).

\section{QA/QC}

\subsection*{Analysis of DNA Yield and Purity}

\NP The Nanodrop...

\NP Open instrument software and click on Nucleic Acid. Login the the software, username: NanoDrop and Password: Nano.

\NP Clean machine by adding distilled water onto the platform and wipe with a Kimwipe before loading anything onto sensor. 

\NP Pipette the X $\mu$L for a blank, using the same solution that was used for the final elution (i.e. Elution Buffer), using a 2 $\mu$L pipettor onto the hydrophobic surface, creating a tiny droplet.

\NP Press Blank.

\NP Clean off surface with a Kimwipe, and repeat the previous step, but replace the elution buffer with each sample to be annalyzed. Be sure the samples are thawed out and tapped to ensure that the sample is homogenized
\NP Load sample and press measure.

\NP Record the 260/280 and 260/230 ratios and the yield (nanograms/L), in your laboratory notebook. They should be around 1.8-2.0 ideally. 


\section{References}

\NP asd fasdf asdf asdf 

\end{document}
